\section{Intellectual property law：Licenses、fair use 与 Lawsuits}

几乎互联网上的原创表达都受 copyright 保护。Copyright 保护 expression，而不是 idea；保护门槛很低，不需要注册即存在，但美国起诉通常需要先注册。使用 copyrighted work 的主要路径是获得 license，或主张 fair use。

\begin{knowledgebox}{课堂提示：License 可以理解为“不起诉的承诺”}
课程先从 intellectual property law 的目标讲起：通过保护作品来激励知识产品的创造。随后给出一个很实用的 license 直觉——licensor 向 licensee 许可某些使用方式，本质上是在约定范围内承诺不诉。Creative Commons 把这种授权标准化；没有许可时，使用者通常只能重新判断 fair use，而不是把“公开可访问”当作默认授权。
\end{knowledgebox}

\begin{importantbox}{术语消化：fair use 四因素}
\begin{enumerate}
    \item 使用目的和性质：教育用途、transformative use 更有利。
    \item 作品性质：事实性、非虚构作品更有利。
    \item 使用数量和实质性：片段比整部作品更有利。
    \item 对原作品市场的影响：若替代原作品市场，更不利。
\end{enumerate}
这些因素需要具体案件具体分析；本笔记只解释课程内容，不构成法律建议。
\end{importantbox}

课程提到的诉讼包括 The New York Times v. OpenAI、作者诉 Anthropic、作者诉 Meta。讲义中的核心 takeaway 是：截至课程材料所述，一些具体案件中训练被认为 fair use，但盗版复制明确违法，且这一领域仍在快速演化。

\begin{warningbox}{训练数据的法律风险不是只看模型是否背诵}
复制数据本身可能已经构成争议；训练是否 transformative、模型是否输出受保护表达、是否影响原市场、是否违反 ToS，都是不同问题。数据管线需要记录来源、许可证、过滤规则和使用理由。
\end{warningbox}

\begin{knowledgebox}{术语消化：训练数据法律/合同术语}
\begin{tabular}{L{0.2\textwidth}L{0.36\textwidth}L{0.34\textwidth}}
\toprule
术语 & 课程中的含义 & 数据工程后果 \\
\midrule
Copyright & 保护原创表达，不保护抽象 idea。门槛很低，网页文字通常也受保护。 & 不能把“公开网页”直接等同于“可训练数据”。 \\
License & 权利人承诺在特定条件下不追究使用。 & 需要保存 license、版本、适用范围和限制。 \\
Creative Commons & 一组允许分发/改编的开放许可证，但不同 CC 条款差异很大。 & 需要区分 CC-BY、CC-BY-SA、NC、ND 等条件。 \\
Fair use & 美国法下四因素平衡判断，是否适用取决于具体事实。 & 不能在数据表里简单标成“fair use = yes”。 \\
Terms of Service & 网站合同条款，可能限制 bot 下载或训练用途。 & 即使版权层面可争辩，合同层面仍可能违规。 \\
\bottomrule
\end{tabular}
\end{knowledgebox}

\begin{importantbox}{如何把法律问题变成数据工程字段}
高质量数据集至少应记录：source URL、crawl time、license 字段、robots.txt/ToS 状态、是否来自 shadow library、是否包含个人信息、处理步骤、过滤器版本、dedup 策略、下游用途假设。没有这些 provenance 字段，后续很难审计风险。
\end{importantbox}

